\documentclass[ecta,nameyear,draft]{econsocart}
\input{preamble}
\begin{document}
\begin{frontmatter}
\title{Neural Bellman Operators}
\runtitle{Neural Bellman Operators}
\begin{aug}
\author[id=au1,addressref={add1}]{\fnms{Qian}~\snm{QI}\ead[label=e1]{qiqian@pku.edu.cn}}
\address[id=add1]{Peking University}
\end{aug}
\begin{abstract}
This paper develops Neural Bellman Operators for constructing, evaluating and improving feasible policies in controlled economies. Own-future continuation construction retains action witnesses and admits policy-identical native and affine--ReLU realizations. We give a directed, acquisition-aware policy-improvement theorem: certified signed action comparisons preserve the incumbent at every state, while a continuous-action lower cover yields a finite-sweep bound for the original Bellman optimum. Coupled residuals construct the comparisons without an unknown policy-value oracle. Compressed pair functions have a verified defect account, and approximate action-null components receive explicit model-error allowances. In the original nonlinear investment economy, full all-date iterations have positive, nonconstant evaluation errors and identify reductions in the witness policies' actual expected costs. Uniform and surplus-adaptive FVI receive the same improvement and evaluation services. Final cross-generator differences are unresolved, not evidence of neural dominance. Complete work-to-cost frontiers, perturbation diagnostics, proofs and the unchanged adverse historical record accompany the results.
\end{abstract}
\begin{keyword}
\kwd{Dynamic programming}\kwd{Neural approximation}\kwd{Policy improvement}
\kwd{Information acquisition}\kwd{Computational economics}
\end{keyword}
\end{frontmatter}
\section{Introduction}\label{sec:introduction53}
A numerical continuation is useful to an economist because it changes a feasible decision. A small prediction error is not the cost of the implemented policy, and a smaller upper bound on that cost is not evidence of a lower cost than a competitor. The distinction becomes consequential when information acquisition, action selection and verification themselves consume resources.

Neural Bellman Operators organize these objects around a common economic recursion. Each date's continuation is constructed from its own fitted future, rather than from the unknown optimal value. A signed policy sandwich evaluates arbitrary fitted candidates. A constructive backend produces an affine--ReLU continuation and a feasible action witness from effective primitive covers and moduli, without presuming successful nonconvex training. Original-owner compilation preserves that action under continuous off-grid queries, and the observation contract specifies the policy that is actually deployed.

The additional problem studied here is how to improve such a constructed policy without losing its guarantee. We certify the signed difference between a proposal's continuation cost and the exact incumbent's continuation cost. An acquired cell is changed only when its entire true-state range is safe. At the same time, a separate continuous-action lower cover measures distance from the original feasible Bellman infimum. All dates are updated simultaneously against a frozen complete policy. The resulting recursion removes the initialization error within the horizon and retains an explicit finite-resolution error sum. It does not replace the economic optimum by a grid optimum.

The central certificate is representation-neutral. A native min-plus envelope and its exactly equivalent ReLU circuit are one constructed method with two encodings. The improvement theorem is also available to conventional FVI. The neural contribution of the construction is therefore stated at the level actually established: an explicit feasible own-future realization and its computational account, not a uniquely neural policy-improvement principle or a generic optimizer-success probability.

Three constructive questions connect the theorem to execution. First, a signed paired-residual recursion evaluates action contrasts without a policy-value oracle; it retains the sign information that an absolute error bound can discard. Second, an on-demand pair tree avoids a square state-pair table, while a residual-defect theorem accounts for low-rank or locally approximated pair functions when their defects are verified. Third, approximate rather than exact action-nullity is charged quantitatively, including errors in the action exposure and innovation marginal laws. The complete work and error costs of these steps are part of the method.

The original nonlinear investment economy supplies a direct experiment. Starting from the original witness and FVI policies, we execute every date of two or three complete improvement passes. State acquisition and policy discontinuities are enclosed on every observation cell, not tested only on random states. The resulting witness policies have identified positive gains relative to their own incumbents. A separately frozen, fresh-runner comparison adds surplus-driven nonuniform FVI and charges its pilot work. The final witness--FVI differences are unresolved. These results establish executed safe improvement, not neural dominance, and do not remove the adverse comparisons for the original policies.

Actual expected cost is the primary empirical endpoint. Independent interval paths evaluate every retained policy under the original continuous initial and innovation laws. Simultaneous intervals permit a common work-to-cost frontier and the full break-even fee account, instead of selecting a convenient fee or subtracting two regret bounds. Whole-cell perturbation diagnostics separately show when omitting a nonzero error allowance accepts a harmful action. They create no extra independent cost observations.

\subsection{Relation to existing methods}
The performance-difference identity and conservative policy improvement are established foundations of approximate dynamic programming; see \citet{kakade2002} and \citet{metelli2021}. We do not claim that a verified nonpositive reference advantage is a new principle. Our additional obligations are a feasible acquired-state implementation, a computable signed enclosure for the entire cell, a lower cover of the original continuous action infimum, and a complete finite-sweep and resource account. Unlike policy mixing around an estimated greedy policy, the executed rule chooses an incumbent-anchored cell action using those enclosures.

State-pair comparisons also have a substantial literature. The bisimulation metrics of \citet{ferns2004} bound value differences by reward and transition discrepancies and support aggregation. Our pair recursion is policy- and critic-specific, finite-horizon and directed: it propagates Bellman residual differences for a fixed incumbent and need not be a metric. It does not require an optimal coupling, but every chosen coupling must have the correct marginals. The compressed-defect bound makes approximation error explicit rather than treating a pair representation as exact.

Potential-based reward shaping preserves optimal policies by a telescoping reward transformation \citep{ng1999}. A critic change in this paper changes no economic reward; action-nullity instead means that the specified action kernels annihilate a continuation component's contrasts. Value equivalence studies models giving the same Bellman updates on selected function and policy classes \citep{grimm2020}. Here the economic kernels are held fixed, and a continuation-error component is removed only after its action contrasts have been verified. These are related cancellation ideas with different quantified objects. Approximate nullity and kernel misspecification require the additional allowances proved below.

Robust baseline improvement explicitly compares a candidate with its baseline under model uncertainty \citep{ghavamzadeh2016}. Our enclosure rule likewise refuses an unverified cost increase, but the primary experiment treats the primitive model as known and certifies numerical and acquisition errors. The separate perturbation catalogue varies explicitly declared kernels; it is not a statistical identification theorem for unknown dynamics.

The own-future constructive backend uses classical Lipschitz extension \citep{mcshane1934}, min-plus approximation \citep{lakshminarayanan2014}, and distance transforms \citep{felzenszwalb2012}. Bellman-inequality approaches also provide value bounds \citep{wang2015}. The present construction retains the feasible original action through those operations and compares the actual returned policies. Its broader controlled-diffusion, recursive-utility, endogenous-preference, temporal-self and game applications remain in the complete NBO development under their stated assumptions. The expectation-specific signed transport and cost-inference theorems are not silently extended to every nonlinear recursion.

\input{sections/core53}
\input{sections/directed_theory53}
\input{sections/resources53}
\input{sections/study53}
\input{sections/decisions53}

\section{Conclusion}\label{sec:conclusion53}
Neural Bellman Operators connect continuation construction to a feasible implemented policy and an economic resource account. The directed certificate adds a concrete way to improve that policy: retain signed reference comparisons, verify them over acquired cells, and independently cover the continuous feasible optimum. Complete all-date iterations in the original investment economy identify own-incumbent gains even with nonzero evaluation errors. The same procedure is available to conventional generators, whose actual costs and full work remain essential comparisons.

The full-sweep bounds do not presume away a fixed information floor, and the pair-tree implementation does not establish dimension-free or non-tensor scaling. The verified compression and perturbation results specify mathematical routes beyond exact nullity and square pair tables. A genuine non-tensor, independently checked dimension experiment and an externally justified economic calibration remain distinct evidence requirements. They are not replaced by a renamed low-dimensional result. The same NBO paper, its original applications and its complete adverse evidence are retained for further mathematical and economic review.

\appendix
\section{Reading and preservation map}\label{app:map53}
This article is the active R53 exposition. The technical supplement supplies the executed identities, full datewise and work records, exact-moment audit, and the earlier technical material. The separate historical article reproduces the complete R52 main manuscript from ordinary preserved sources, including every theorem, resource-allocation result, information decision and adverse table. The complete development and complete proof supplement retain the full original bodies and add the present results. This organization shortens the active argument by relocating, not deleting, the cumulative exposition.

The revision contains the pre-execution protocols, immutable scientific sources and checkpoints, all raw interval and policy records, generated tables and the point-by-point response. Source hashes, exact tests, numerical audits, complete document builds and an independent clean Git-archive rebuild are bound in the release records. Their purpose is to supply a reproducible object for scrutiny; they are not an editorial acceptance decision or a substitute for independent mathematical review.
\bibliographystyle{ecta-fullname}
\bibliography{references}
\end{document}
